\documentclass{article}
\usepackage[a4paper, margin=2.5cm, headheight=13pt]{geometry}
\usepackage[french]{babel}
\usepackage{fancyhdr}
\usepackage{tikz}
\usetikzlibrary{calc}
\usepackage{graphicx}

% Définir la pagination et.c.
\pagestyle{fancy}
\fancyhf{}
\lhead{MICRO-315 Miniprojet} 
\rhead{Les aventures de Jean-Bot}
\fancyfoot
{
    \begin{tikzpicture}[remember picture,overlay]
            \node[] (chat) at ($(current page.south)+(0,1.5)$) {\includegraphics[width=1cm]{images/amour.eps}};
            \node[] () at ($(chat)+(0,0.04)$) {\thepage};
            \node[] () at ($(chat)-(1,0)$) {\includegraphics[width=1cm]{images/kawwai.eps}};
            \node[] () at ($(chat)+(1,0)$) {\scalebox{-1}[1]{\includegraphics[width=1cm]{images/kawwai.eps}}};
    \end{tikzpicture}
}

\begin{document}
    \begin{titlepage}
        \vspace*{\fill}
        \begin{center}
            \textbf{\large MICRO-315 Miniprojet}
            \vspace{2.5cm}
            
            \begin{tikzpicture}[remember picture,overlay]
                \node (titre) at (0,0)
                    {
                        \begin{tabular}{c}
                            \hline
                            \\
                            \huge\it Les aventures incroyables de \\
                            \\
                            \Huge\bf Jean-Bot I Baptiste\\
                            \\
                            \huge\it Le raccourceur des circuits\\
                            \\
                            \hline
                        \end{tabular}
                    };
                \node[left]  at (titre.west) {\includegraphics[width=5cm]{images/kawwai.eps}};
                \node[right] at (titre.east) {\scalebox{-1}[1]{\includegraphics[width=5cm]{images/kawwai.eps}}};
            \end{tikzpicture}

            \vspace{3cm} 
            \textit{\large Un projet de robotique de Julian Donevsky et Olof Lissmats}
        \end{center}
        \vspace{5cm}
        \vspace*{\fill}
    \end{titlepage}

    \thispagestyle{empty}
    \tableofcontents

    \newpage
    \setcounter{page}{1}

    \section{Introduction}
    Dans ce projet on a aidé notre cher ami Jean-Bot de développer ces compétences de danse, chanson et géographie.
\end{document}
